\documentclass{article}
\usepackage{natbib}

\title{Deep Learning Approaches for Protein Structure Prediction}
\author{Jane Doe}
\date{2026}

\begin{document}
\maketitle

\begin{abstract}
Recent advances in deep learning have revolutionized protein structure
prediction. Methods based on neural networks, particularly transformer
architectures, have achieved remarkable accuracy in predicting
three-dimensional protein structures from amino acid sequences. This paper
reviews the key developments in this field, including AlphaFold and
related approaches.
\end{abstract}

\section{Introduction}
Protein structure prediction has been a grand challenge in computational
biology for decades. The ability to predict the three-dimensional structure
of a protein from its amino acid sequence has profound implications for
drug discovery and understanding biological function.

\section{Methods}
Transformer-based architectures have shown remarkable performance in
capturing long-range dependencies in protein sequences. Attention
mechanisms allow the model to focus on evolutionarily related residues
that co-evolve across species.

Multiple sequence alignments provide crucial evolutionary information
that constrains the space of possible structures. These alignments are
processed through specialized neural network layers.

\section{Results}
Our experiments demonstrate that combining attention mechanisms with
evolutionary features significantly improves prediction accuracy on
CASP benchmarks.

\section{Conclusion}
Deep learning has transformed protein structure prediction, enabling
accurate modeling at scale.

\bibliographystyle{plain}
\bibliography{references}

\end{document}
